\chapter{Split-host физичка архитектура и MotionSDK}
\label{chap:hardware}

\section{Одделување на физичкиот пат}

Физичката интеграција не е продолжение на Gazebo launch-от. Таа користи
посебен catkin workspace на perception computer-от и не ги менува vendor
\code{lite\_cog}, \code{qnx2ros}, \code{jy\_exe} или нивната конфигурација.
Како што е прикажано на \cref{fig:split-host}, reasoning и OpenAI остануваат на
развојниот компјутер, додека телеметријата, safety gates, trajectory generation
и единствениот MotionSDK owner се на роботот.

\begin{figure}[htbp]
\centering
\resizebox{\textwidth}{!}{\begin{tikzpicture}[node distance=18mm and 25mm]
  \node[process,text width=48mm] (dev) {
    \textbf{Development computer}\\
    chat + EmotionEngine\\
    OpenAI или deterministic reply\\
    hardware uplink};
  \node[physical,right=of dev,text width=48mm] (perception) {
    \textbf{Perception computer}\\
    state validation + telemetry\\
    contact/safety gates\\
    sole MotionSDK runner};

  \node[physical,below=of perception,text width=48mm] (motion) {
    \textbf{Motion host}\\
    vendor \code{jy\_exe}\\
    state + joint feedback};
  \node[physical,left=of motion,text width=48mm] (robot) {
    \textbf{Physical Lite3}\\
    12 joints + IMU\\
    Retroid / STOP evidence};

  \draw[flow] (dev) --
    node[above,font=\scriptsize]{SSH loopback: NDJSON 1.0}
    (perception);
  \draw[flow,<->] (perception) --
    node[right,font=\scriptsize,align=left]{MotionSDK UDP\\exclusive owner}
    (motion);
  \draw[flow,<->] (motion) --
    node[below,font=\scriptsize]{vendor actuation / feedback path}
    (robot);
  \draw[dashedflow] (robot.north east) --
    node[above,sloped,font=\scriptsize,fill=white,inner sep=1.5pt]{validated 0x0901 / 0x0906 telemetry}
    (perception.south west);
\end{tikzpicture}
}
\caption{Split-host физичка архитектура и насоки на податоците.}
\label{fig:split-host}
\end{figure}

Development-side \code{hardware\_brain.launch} стартува само chat adapter,
emotion adapter и hardware uplink. Uplink-от праќа NDJSON schema 1.0 кон
\code{127.0.0.1:8767} преку експлицитен SSH local forward. Frame е до 2 KiB,
има random session и strictly increasing transport sequence, а внатре ја носи
неизменета emotion-state 1.1. На perception computer-от
\code{emotion\_receiver.py} отфрла replay и malformed state.

Симулацискиот и физичкиот пат се споредени во \cref{fig:path-comparison}.
Нивна единствена заедничка actuator-near граница е валидираната емоционална
состојба. Затоа \code{/emotion\_bot/joy\_out} и симулацискиот
\code{safety\_bridge.py} не се користат на hardware.

\begin{figure}[htbp]
\centering
\resizebox{\textwidth}{!}{\begin{tikzpicture}[node distance=10mm and 9mm]
  \node[box,text width=29mm] (state) {Emotion state 1.1\\заеднички contract};
  \node[process,above right=4mm and 13mm of state,text width=28mm] (map) {Twist mapper};
  \node[safety,right=of map,text width=28mm] (safe) {simulation safety};
  \node[process,right=of safe,text width=27mm] (joy) {Joy transport};
  \node[process,right=of joy,text width=25mm] (gz) {Gazebo};
  \node[font=\scriptsize\bfseries,above=2mm of map] {Симулација};
  \draw[flow] (state.north east) -- (map.west);
  \draw[flow] (map) -- (safe);
  \draw[flow] (safe) -- (joy);
  \draw[flow] (joy) -- (gz);

  \node[physical,below right=4mm and 13mm of state,text width=28mm] (resolve) {profile resolver};
  \node[safety,right=of resolve,text width=28mm] (gates) {robot/contact gates};
  \node[physical,right=of gates,text width=27mm] (sdk) {12-joint SDK};
  \node[physical,right=of sdk,text width=25mm] (hw) {Physical Lite3};
  \node[font=\scriptsize\bfseries,below=2mm of resolve] {Физички робот};
  \draw[flow] (state.south east) -- (resolve.west);
  \draw[flow] (resolve) -- (gates);
  \draw[flow] (gates) -- (sdk);
  \draw[flow] (sdk) -- (hw);
\end{tikzpicture}
}
\caption{Споредба на симулацискиот и физичкиот контролен пат.}
\label{fig:path-comparison}
\end{figure}

\section{Телеметрија, STOP и ексклузивна сопственост}

\code{telemetry\_tap.py} пасивно ги набљудува пакетите од motion host кон UDP
43897. Проверениот Deeprcs 2.0.153 layout содржи 368-byte high-rate
\code{0x0906} body: tick, IMU, 12 joints и contact block. Независниот
\code{0x0901} state body има 200 bytes и се користи за basic state, gait/motion,
battery и error flags. Само валидирани записи се објавуваат во sequence-locked
shared memory.

Motion-host STOP observer работи без ROS, пасивно го набљудува Retroid traffic
и праќа HMAC-authenticated status до perception computer. Само tap-от и STOP
observer-от добиваат \code{CAP\_NET\_RAW}; ROS core и runner остануваат
непривилегирани. Stale relay ја отстранува arming дозволата.

Во една operator-started сесија постои точно еден
\code{motion\_sdk\_expression\_runner} на 1 kHz. Ownership marker и
independent crash watchdog спречуваат втор sender. Retroid diagnostic bridge,
legacy posture/action nodes и официјалниот runner се меѓусебно исклучиви.
Expression status се враќа преку shared memory и read-only JSON schema 1.0 на
\code{/emotion\_bot/hardware/expression\_status}.

\section{Од high-level packets до директни joints}

Раната истрага утврдила дека full Lite3 Android app испраќа 12-byte
little-endian \code{<III>} commands кон motion host \code{192.168.2.1:43893}.
Тоа не е generic Retroid port-12121 protocol. State-aware Stand/Sit е toggle и
не смее слепо да се повторува. Потребни се heartbeats и fresh
\code{robot\_basic\_state}.

Три height-only обиди низ
\code{/emotion\_bot/hardware/posture\_intent -> /simple\_cmd -> ros2qnx}
не произвеле видливо движење. По serialization fix, nonzero height packets и
точна zero recovery биле видливи на ROS и perception wire, но joint span бил
само 0.00099 rad и операторот повторно пријавил нулто значајно движење. Auto
mode neutral breathing исто дал само 0.00084 rad span. Ова докажало command
publication и wire emission, но не \code{jy\_exe} acceptance.

Подоцнежен direct Retroid-compatible high-level test ја репродуцирал exact
sequence: четири heartbeats на 2 Hz, Move, 2.0 s, Pose, 1.5 s и потоа 50 Hz
pairs од opaque yaw code \code{0x21010135}=32768 пред height
\code{0x21010102}. Фиксна вредност $-10000$ дала 0.2484894 rad joint span, со
state 6 и нулти errors. Тоа е доказ за bounded height actuation, но periodic
profile не бил докажан како дишење; Retroid diagnostic останува одделен од
официјалниот пат.

\section{MotionSDK takeover и официјална секвенца}

Официјалниот MotionSDK документира UDP communication, 12-joint order и
контролно преземање. Критичен наод е
дека \code{RobotStateInit()} не ја чита и не ја задржува моменталната поза. Тој
ги reset-ира joints кон zero и презема control right. Затоа current-pose/ground
takeover е небезбеден.

Еден ground measured-hold обид предизвикал backward stepping и joint-error
watchdog за приближно 125 ms. Друг „мал offset“ takeover ги засилил барањата од
ред $10^{-5}$ rad во spans до 0.327/0.709/0.677 rad по joint families. Обидот е
прекинат и \code{takeover\_transition\_commissioned=false} останал hard gate.
Овие резултати се причина официјалниот пат да започнува само од exact basic
state 1 и да ја користи vendor zero-to-stand секвенцата:

\begin{equation}
\begin{aligned}
\texttt{RobotStateInit} &\rightarrow \texttt{PreStandUp}
\rightarrow \texttt{StandUp} \\
&\rightarrow \text{neutral hold}
\rightarrow \text{expression} \rightarrow \texttt{ControlGet(ROBOT)}.
\end{aligned}
\label{eq:official-sequence}
\end{equation}

Секвенцата се извршува еднаш по сесија. Категориите се менуваат под ист owner;
не се повторува stand и не се reacquire-ира SDK.

\section{Feedback watchdog и hard interlocks}

High-rate feedback age над 100 ms го замрзнува trajectory time и ја задржува
последната валидирана position со нулти desired velocity и torque. Продолжување
е дозволено по 20 свежи distinct frames со age до 20 ms. Dead feedback на
250 ms, non-finite data, tracking violation, state 8 или nonzero error веднаш
водат кон no-retry release. Attitude bound е $\pm10^{\circ}$, а battery floor
е експлицитно конфигуриран за конкретна сесија.

Важно е дека software STOP не е обична timeout recovery постапка. STOP,
invalid state, tracking fault и други hard faults го заобиколуваат
разговорниот 1.5+0.35 s transition contract и го активираат проверениот
hold/release. Системот никогаш автоматски не reacquire-ира по release.

\section{IK и torque-derived contact}

Физичкиот ROS bridge нема mode-valid foot-contact topic, а MotionSDK contact
array бил all-zero под SDK ownership. Затоа отворен foot lift не бил прифатен.
Контактот се проценува од joint position и torque со

\begin{equation}
  J(q)^{T}F=\tau, \qquad F=(J(q)^{T})^{-1}\tau,
  \label{eq:contact-estimator}
\end{equation}

при што vertical load е $\max(0,-F_z)$. Estimator-от отфрла non-finite,
singular и implausible solutions, користи distinct ticks, filtering и
consecutive-frame hysteresis. Recorded standing fixture дала
23.44/24.34/31.23/35.73 N, вкупно 114.74 N, наспроти приближно 116.15 N за
11.84 kg.

За paw gesture се бара: valid four-foot baseline, target unload, најмалку два
strong non-target supports над 5 N, најмалку 70 N aggregate non-target load,
bounded IK lower, target landing и повторно четири supports. Contact
confirmation не е замена за state, attitude, tracking, battery, STOP или
feedback gates; сите мора истовремено да важат.
